\begin{table}[htbp]
\caption{Diagnóstico del modelo logístico}
\label{tab:tabla_24_diagnostico_logistica}
\begin{tabular}{lll}
\toprule
aspecto & valor & veredicto \\
\midrule
Eventos por variable (EPV) & 104.2 & Holgado \\
Discriminación (AUC) & 0.667 (0.655–0.679) & discriminación aceptable \\
Bondad de ajuste (Hosmer-Lemeshow) & χ² = 21.42, gl = 8, p = 0.0061 & Rechaza el ajuste, pero con N grande la prueba es hipersensible: juzgar la calibración por la curva y el Brier score \\
Calibración (Brier) & 0.0569 & Brier 0.0569 frente a 0.0582 del modelo nulo (mejora del 2.2 \%) \\
Linealidad del logit en edad (Box-Tidwell) & p = 0.0000 & El logit NO es lineal: usar splines o categorizar \\
Sensibilidad (corte en la prevalencia) & 0.672 & Umbral 0.0621; el modelo está calibrado, así que el corte en la prevalencia es interpretable \\
Especificidad (mismo corte) & 0.564 & Compromiso propio de un tamizaje: capta más casos a costa de más falsos positivos \\
AUC-PR & 0.1169 & Frente a 0.0621 del azar: 1.9 veces mejor \\
Exactitud del modelo trivial & 0.938 & Referencia obligada: predecir 'nadie' acierta el 93 \% \\
\bottomrule
\end{tabular}
\end{table}
